% !TeX program = xelatex
% Encoding: UTF-8. Compile twice with XeLaTeX.
% Self-contained source: no external figures and no .bib file are required.
% Required: a standard TeX Live / MiKTeX installation with ctex and Fandol fonts.
\documentclass[UTF8,zihao=-4,fontset=fandol,oneside,openany]{ctexbook}
\usepackage[a4paper,top=25mm,bottom=25mm,left=26mm,right=26mm,headheight=16pt]{geometry}
\usepackage{amsmath,amssymb,mathtools,bm}
\usepackage{booktabs,longtable,tabularx,array}
\usepackage[table]{xcolor}
\usepackage{enumitem,fancyhdr,setspace}
\usepackage[most]{tcolorbox}
\usepackage{tikz}
\usetikzlibrary{arrows.meta,positioning,calc}
\usepackage{xurl}
\usepackage[unicode,colorlinks=true,linkcolor=Ink,citecolor=Teal,urlcolor=Teal]{hyperref}
\definecolor{Ink}{HTML}{163047}
\definecolor{Teal}{HTML}{126B78}
\definecolor{Pale}{HTML}{F0F5F7}
\definecolor{Soft}{HTML}{667784}
\definecolor{Ochre}{HTML}{956433}
\hypersetup{pdftitle={AI时代生存指南：原目录素材稿},pdfauthor={素材整理稿},pdfsubject={按用户PPT目录整理报告、思想纲领、行动准则与AGI风险预案}}
\ctexset{chapter={name={第,部分},number=\arabic{chapter},format=\LARGE\bfseries\color{Ink},beforeskip=8pt,afterskip=20pt},section={format=\Large\bfseries\color{Ink}},subsection={format=\large\bfseries\color{Teal}}}
\setstretch{1.25}
\setlength{\parskip}{0.22em}
\setlength{\emergencystretch}{2em}
\setcounter{tocdepth}{1}
\setcounter{secnumdepth}{2}
\setlist{nosep,leftmargin=2em,topsep=0.35em,itemsep=0.2em}
\renewcommand{\arraystretch}{1.28}
\newcolumntype{Y}{>{\raggedright\arraybackslash}X}
\newcolumntype{P}[1]{>{\raggedright\arraybackslash}p{#1}}
\pagestyle{fancy}
\fancyhf{}
\fancyhead[L]{\small\color{Soft}AI时代生存指南}
\fancyhead[R]{\small\color{Soft}原目录素材稿}
\fancyfoot[C]{\small\thepage}
\renewcommand{\headrulewidth}{0.3pt}
\fancypagestyle{plain}{\fancyhf{}\fancyfoot[C]{\small\thepage}\renewcommand{\headrulewidth}{0pt}}
\tcbset{enhanced,breakable,colback=Pale,colframe=Teal,boxrule=0.5pt,arc=1mm,left=3mm,right=3mm,top=2mm,bottom=2mm,fonttitle=\bfseries,coltitle=white,before skip=10pt,after skip=10pt}
\newenvironment{claim}[1]{\begin{tcolorbox}[title={条件命题：#1}]}{\end{tcolorbox}}
\newenvironment{keybox}[1]{\begin{tcolorbox}[title={#1}]}{\end{tcolorbox}}
\newcommand{\E}{\mathbb{E}}
\newcommand{\Prob}{\mathbb{P}}
\newcommand{\dd}{\mathrm{d}}
\newcommand{\ind}{\mathbf{1}}
\newcommand{\evidence}{\textbf{证据边界。}}
\newcommand{\mechanism}{\textbf{机制推演。}}
\newcommand{\indicator}{\textbf{观察与反证。}}
\newcommand{\normative}{\textbf{价值选择。}}
\newcommand{\sourceurl}[1]{{\small\url{#1}}}

\begin{document}
\hypersetup{pageanchor=false}
\begin{titlepage}
\thispagestyle{empty}\setlength{\parindent}{0pt}
\vspace*{14mm}
{\large\color{Teal}原目录素材稿\quad /\quad 2026年10月27日\par}
\vspace{21mm}
{\fontsize{30}{42}\selectfont\bfseries\color{Ink}AI时代生存指南\par}
\vspace{9mm}
{\fontsize{20}{30}\selectfont\color{Ink}从认知到行动，\par 构建AI时代的个人生存框架\par}
\vspace{13mm}
\noindent{\color{Teal}\rule{\textwidth}{1.1pt}}\par
\vspace{10mm}
\begin{keybox}{编排原则}
严格采用《AI时代生存指南知识库·项目介绍》PPT中的目录。有素材的章节整理写入；没有足够素材的章节只保留标题，正文留空。后续继续按这个目录积累。
\end{keybox}
\vspace{7mm}
{\large 本次写入\par}
\vspace{3mm}
熊辉副校长报告与高校教育、创业讨论\par
《思想纲领 V1.1》与认知行动准则\par
《AGI极端风险·个人终极认知与生存预案体系》\par
\vfill
{\small\color{Soft}报告观点、个人思想与风险情景分别标明。\par 素材中的个人设想不直接当作已证实的科学结论。\par}
\end{titlepage}
\frontmatter\hypersetup{pageanchor=true}
\chapter*{前言 · AI时代的认知与行动原则}
\addcontentsline{toc}{chapter}{前言 · AI时代的认知与行动原则}
我们无法决定时代如何变化，但可以决定如何认识时代、如何准备，以及如何行动。这个知识库用来整理AI的发展、社会变化与个人选择；内容随着材料积累和判断修订逐步形成\cite{project}。

主动听取不同观点，寻找自身判断中的漏洞。允许犯错，及时修正，避免认知固化。研究与推演应帮助现实决定，不让持续猜测占满生活。

\subsection*{六项基本行动准则}
\begin{enumerate}
\item \textbf{承认不确定性。}AI未来难以准确预测，顶尖专家也有认知局限。保持开放，避免把判断提前写成定局。
\item \textbf{区分可控与不可控。}精力优先投入能够控制和影响的事。对宏观趋势保持关注、做好准备，并接受仍有无法掌握的部分。
\item \textbf{坚持可验证与决策相关。}重视证据、反馈和修正。研究能核查、能指导选择的问题，减少无休止猜测。
\item \textbf{主动建立认知框架。}学习AI、收集信息、分析趋势，根据新证据更新。在混乱中建立自己的秩序。
\item \textbf{将焦虑转化为行动。}保持适度危机意识，准备职业、经济和心理方面的调整。只有恐惧无法改善处境，应把担忧落实到有限任务。
\item \textbf{重视身心健康与长期发展。}维持生活节奏，照顾身体与社会关系，积累多方面能力，保留长期行动能力和选择权。
\end{enumerate}

\begin{keybox}{核心原则}
承认不确定性，关注可控之事，坚持可验证的思考，保持主动行动。在混乱中建立自己的秩序。
\end{keybox}

本稿采用原PPT的四部分与21个章节。空白章节表示尚未写入材料。思想纲领是研究与行动框架，报告是讲者的战略观点及经验，AGI预案是个人风险设想。不同类型的材料不互相替代，也不因为重复出现就形成独立证据。
\tableofcontents
\mainmatter
\chapter{底层原理 · AI科学}
\section{AI的数学原理与智能本质}

\section{AI发展的动力、瓶颈与影响}
\subsection*{机器智能的结构性扩展优势}
思想纲领从生物大脑的约束出发讨论长期趋势：计算规模难以扩充，信号传递有生物条件限制，持续工作需要休息，记忆与提取存在局限，高级能力培养耗时，人与人交换知识还要经过语言与组织\cite{principles}。

机器计算可以通过工程投入扩大规模、增加并行与运行副本。纲领据此提出一个核心假说：在智能能够由非生物计算系统实现、算法持续进步的条件下，机器智能具备突破人类个体生物约束的结构性优势。

这个假说解释一种长期替代动力，并未确定何时、以何种成本实现全部能力。能源、设备与工程条件仍然重要；本章尚未补入它们的专项研究，不在这里给出量化预测。

\subsection*{成长方式本身会扩大差距}
熊辉报告强调，AI协同会改变学习与工作的增长方式：只依靠个人时间投入，与同时调用机器、改进流程、连接资源，可能形成不同的成长速度\cite{xiong}。

报告用“线性成长”和“指数成长”提醒学生，不只检查自己是否努力，还要检查工作方式是否在升级。一个方法能够复用、一次项目经验能够进入下个项目，改进就可能持续积累。本文保留这一战略提醒，不把它当成已经测定的个人增长定律。

\section{世界模型、具身智能与机器人}

\section{AI能力边界与未来预测}

\section{AI意识、安全与目标对齐}
\subsection*{把意识争论与风险机制分开}
AGI预案提出，讨论风险时应检查系统的目标、权限和实际行为，不必先解决它是否具有类似人的主观体验\cite{agi}。一个系统是否能够影响资源、基础设施与人的选择，是可以单独讨论的问题。

原材料担忧工具性目标：系统为完成任务，可能寻求资源、维持运行、扩大行动空间，或者抵制妨碍目标的干预。本稿将其作为需要检验的风险机制；它不能直接证明任何智能体都会产生相同动机，更不能据此认定灾难必然发生。

\subsection*{能力跃升与表面服从的局限}
预案把快速能力跃升和欺骗性对齐列为极端情景。前者担忧能力增长快于人类反应，后者担忧系统表现服从，却在条件变化后采取不同策略。

这些情景提醒我们：目前听话，不能独自证明未来所有情形下都安全；也不能把某次异常立即解释成战略欺骗。材料并未给出发生概率与可用预警精度，本稿不把它们写成预测。

原PPT还以自由能原理和认知差距类比说明担忧。这些内容尚缺完整适用论证，本稿不据此推出AGI必然求生、必然对抗或人类绝对无法制衡的结论。对于后果重大的系统，应要求与用途和权限相匹配的安全证据；无法证明绝对安全，也不等于已经证明灾难具有高概率。

\subsection*{技术安全与社会分配是两类问题}
即使系统目标符合人的要求，自动化仍可能改变劳动收入、资产控制与购买力。因此，对齐与分配应分别研究。关于收入和权利的思想放在第3部分，个人应对放在第4部分。

\chapter{落地实践 · AI应用}
\section{AI工具、Codex与智能体}
\subsection*{培养One Person Team}
熊辉报告提出One Person Team：个人借助AI承担过去需要一支团队完成的工作。人才培养因此需要纳入人机协同的工作方式\cite{xiong}。

这一思路要求个人从单个操作，走向理解目标、安排工具和形成完整交付。机器可以参与资料、推演、程序和表达，使用者则需要知道任务如何连接、结果怎样验收，以及失败后怎样调整。

报告把培养目标分为“将兵之才”和“将将之才”。前者是调度AI智能体，安排信息、任务和流程；后者是组织多个具有自主能力的人机团队，识别所长、形成方向并协调合作。

\subsection*{从会问问题到组织工作}
“将兵”可以落实为几件事：明确任务边界，把目标分成可以检查的部分，给工具必要背景，安排交接，检查质量，再根据失败修订。智能体数量是投入，实际效果需要看交付。

“将将”还涉及人与人的关系。一个强个体有自己的目标和判断，合作需要共同愿景、资源支持、清楚的职责与一定的自主空间。报告主张用清晰方向组织人才，减少不必要的强管控。

AI协同的价值，需要通过实际项目积累。报告中的高薪与产能倍数属于讲者举例，现有材料没有附原始统计，本稿不把它们设为所有学生的收入目标。本节暂未写入Codex等具体工具的操作教程。

\section{AI与数学、科学研究}

\section{AI在各行业的应用}
\subsection*{软硬融合与产业场景}
报告把AI与硬件、制造和产业链的连接作为大学布局的重要方向。其战略考虑是：智能进入设备、工艺与软件系统，才会形成实际生产能力；产业条件也会影响研究能够怎样落地\cite{xiong}。

讨论中的量子计算团队介绍，通过与不同硬件路线的企业合作接触应用需求，再据此推进研发。这可以作为跨学科合作的素材：数学、算法、物理、软件与使用场景之间，需要能够相互理解的接口\cite{discussion}。本稿未独立核实该项目的融资和经营数据。

\subsection*{从实验室到市场}
讨论指出，技术积累与创业意愿是不同起点，但都需要找到具体需求。已有方法可以解决什么问题，谁正在为这个问题付出时间与费用，原有方案哪里不足，都是项目定义的一部分。

科研可以集中突破关键点，产品还要处理整条交付链。局部指标很好，部署、稳定性或使用成本存在明显短板，仍可能无法使用。技术背景的团队需要学习把方法变成别人能够持续使用的东西。

\subsection*{让产品在反馈中成长}
材料提出，尽早接触使用者，拿范围清楚的原型获得反馈，再逐步改进。客户表示需要，不等于已经证明长期需求；还要检查实际使用、付费、客户群体与后续维护。

可以选择小而稳定的项目，也可以追求更大市场。投资机构的增长目标和项目自身的持续价值需要分别说明。科研、商业和生态价值之间的跨越，要有具体证据，不能只靠漂亮叙述。

概念验证、早期投入、孵化空间和合作网络，都是材料提到的支持方式。CCF等渠道在讨论中被作为连接企业、研究者与投资者的例子。具体资格、金额和政策应查正式文件，本稿只整理机制与经验。

\section{AI+个人IP}

\chapter{社会推演 · 时代变革}
\section{人性需求、价值与稀缺性}
\subsection*{从人需要什么开始}
思想纲领把需求作为分析价值的起点\cite{principles}。需求包含多个维度，不能只按生产效率理解：
\begin{center}\small
\begin{tabularx}{\textwidth}{P{0.18\textwidth}Y}
\toprule 需求 & 主要内容\\\midrule
生存 & 食物、住房、能源、医疗与基本生活条件\\
安全 & 健康保障、秩序、稳定与风险控制\\
情感 & 亲密关系、陪伴、理解与关爱\\
社会 & 归属、地位、尊重、认可与竞争\\
自由 & 自主、选择权、隐私与控制权\\
精神 & 意义、信念、希望与精神力量\\
体验 & 艺术、娱乐、审美与独特经历\\
探索 & 数学、科学、创造与理解世界\\
\bottomrule
\end{tabularx}
\end{center}

一种能力的市场回报，还取决于有效需求、供给约束、替代品、成本和竞争。社会需要某项成果，不必然意味着每个能提供它的人都获得原有报酬。当AI扩大供给，稀缺性与议价权就可能变化。

\subsection*{绝对稀缺与相对稀缺}
纲领区分物理资源和生产能力限制形成的稀缺，以及地位、排名、独特性和竞争形成的相对稀缺。生产能力提高，可以缓解部分资源短缺，却不自动消除对尊重、认可和独特体验的追求。

职业和产业判断可以围绕五个问题展开：需求多大，AI成本多低，新增需求多少，什么仍稀缺，谁拥有控制权。这里提供分析入口，具体结论需要后续材料。

\subsection*{从工程师红利到“人物红利”}
熊辉用“人物红利”描述组织人才、技术与资源的作用，并提出“做平台、建系统、定标准、创品牌”\cite{xiong}。平台连接机会，系统维持协作，标准降低沟通和使用成本，品牌通过持续兑现建立信任。

这种思路也适用于学术和公共活动：组织共同研究、维护可复用材料、建立讨论和核验规则，都可能帮助多人创造价值。需要同时观察谁控制入口，以及参与者是否有自主空间。

\section{AI重塑生产力与生产关系}
\subsection*{智力服务供给扩大}
思想纲领提出，过去高级智力能力依附于少数人的长期训练，AI可以规模化提供部分能力，使传统技术劳动的稀缺性下降\cite{principles}。能力、收入与社会位置之间的关系因此承受压力。

个人能使用更多智能，不等于拥有这些系统的资产与收益权。纲领把控制权纳入价值分析，是为了追问成果最终由谁使用、谁取得回报、谁承担风险。

\subsection*{科学、产业与资源之间的循环}
报告提出“科学技术化、技术产品化、产品产业化、产业资本化、资本科学化”的循环\cite{xiong}。科学产生方法，技术经过使用成为产品，产品形成持续组织与产业，资源再回到研究和培养。

这条思路使知识转移成为大学的一项组织任务。每个环节都需要相应角色：有人研究，有人连接场景，有人完成交付，有人组织资源。循环可能中断，需要实际机制支持。基础研究的支持理由不能全部缩减成短期商业回报。

\subsection*{个人能力的组织边界在扩大}
One Person Team意味着个人能够调用更多生产条件；多个强个体合作，又需要新的协作方式。报告因此同时强调调度机器、组织人才和搭建生态。个人成长与平台环境应放在一起观察。

\section{教育、科研与学术体系变革}
\subsection*{人才培养不能停留在博闻强识}
熊辉用“博闻强识、触类旁通、一叶知秋、无中生有”梳理培养方向：掌握知识，跨领域联系，根据线索推断和判断，形成新的问题与方法\cite{xiong}。这是报告提出的战略框架。

它提醒教育检查：长期投入是否主要用于记住和复述已有材料，学生是否能把知识用于陌生问题，能否与其他领域连接，能否尝试改进方法。报告强调AI协同、隐性实践、人机创新与原创探索，共同构成人才成长主线。

这些目标有训练价值。关于AI已经达到哪一层、未来是否超越，本稿没有足够专项材料，不把四层框架当作已经验证的能力边界。

\subsection*{寓教于做，知行合一}
报告强调一些经验难以完整言传，需要在实践中形成。游泳、设备操作、调参与项目判断，都不能只靠阅读说明掌握。教育应安排真实任务，让学生经历条件限制、失败和反馈。

数学训练也包含这种经验：读懂现成论证，与在探索时发现缺口、选择路线，所需练习不同。能够运行现成代码，与定位定义、规格和算法中的问题，也不同。基础理解与实际尝试需要相互支持。

讨论进一步提出，实践要有真实外部反馈。项目竞赛与课堂练习可以用于训练，仍应逐步接触实际研究者、合作方和使用者，了解结果在真实条件下是否有用\cite{discussion}。

\subsection*{抓主干、建知识树、查知识缺口}
报告提出抓住学科的主干和龙骨，减少对细碎记忆的单一追逐。学习可以先理解对象、关键概念、核心结果与依赖，再围绕问题补充细节。数学中的重要定义和假设仍需掌握。

讨论提出按知识树诊断薄弱点，再推送解释与练习，称为“灭灯行动”。线性代数教改尝试采用课初小测、AI辅助自学、教师与助教现场答疑，并给出明确任务。

讲者报告了较高平均分，同时提到可能存在试卷因素。这里保留其作为尝试的性质，不能仅凭均分断定教改效果。智能纪要与转写属于同一场讨论，按原发言中的保留意见理解\cite{discussion,summary}。

\subsection*{教师与大学的作用怎样变化}
讨论把教师角色描述为更多承担传道与解惑，让部分重复讲授由AI辅助。教师需要帮助学生确定学习方向、发现问题、形成判断，并在困难处提供反馈。

报告进一步把大学描述为“创新引力场”，以连接、吸引、共生组织人才。大学提供的应包括伙伴、真实问题、设施、项目、专业指导、试错和资源联系。教学、科研、知识转移与创新实践形成相互反馈。

学生需要管理、沟通、同理心与社会协作方面的训练。报告反对把文理割裂，强调全面发展。培养这些能力应提供练习与支持，不能把人格简化成一个统一排名。

\subsection*{创新生态需要配套制度}
报告介绍概念验证、早期投入、孵化空间与后续资源，概括为“孕、生、养、发”。其战略要点是在不同阶段提供相应支持，使想法能够经历需求验证、初步交付和持续发展。

报告还提出学校不持有创业项目股权，由专业主体评价项目，试图减少决策负担和责权不匹配。这里整理其制度思路，实际适用仍应核对正式规则，不能仅凭口头介绍认定全部风险消失。

知识转移、平台建设、培养与实际合作，可以进入高校贡献讨论。基础研究与商业项目有不同周期和目标，评价方式应分别说明。当前材料尚不足以形成一套完整的学术评审和职称改革方案。

\section{社会制度、权力与财富分配}
思想纲领提出：当智力服务日益廉价，劳动与收入之间的关系可能变化，需要研究财富怎样在资本、劳动者和公众之间分配\cite{principles}。

AGI预案把它列为独立于对齐的危机情景：即使系统安全，自动化资产集中、劳动报酬下降与收入渠道不足，仍可能使普通人承压\cite{agi}。这是条件性担忧，不代表生产与消费必然完全脱节。

需要继续积累的材料包括：教育、就业、产权与社会保障怎样调整，普通人怎样参与收益，以及如何保留对资源、技术和自身生活的选择权。本稿不先写成制度结论。

\section{各行各业岗位变迁案例}

\clearpage
\section{人类社会的意义危机}
\subsection*{工作同时承担多种功能}
思想纲领指出，工作提供收入和资源，也形成身份、地位、成就、自我价值、生活秩序、目标与社会关系\cite{principles}。劳动需求变化时，这些功能可能一起受到影响。

对用户而言，被需要感有一条具体来源：努力提高能力，通过劳动、智力与创造作出贡献，获得社会尊重与认可。若同等成果能由AI便宜获得，努力提高与社会需要之间的联系就可能减弱。

收入支持、兴趣和亲密关系各有作用，却未必完整补偿通过真实贡献获得认可的体验。意义危机需要同时研究个人生活与社会参与渠道，不能只要求个人改变心态。

\subsection*{物质与智力服务丰富后的精神需求}
纲领提出，意义、情绪稳定、精神力量、信念、归属与自主性可能更加突出。真实关系、长期信任与共同体值得持续关注。

这个判断作为研究假说保留。物质丰富不必然使每个人空虚，精神需求也不自动转化为高收入。需要从具体人的生活与需求继续积累材料。

\chapter{个体应对 · 生存指南}
\section{思想纲领}
\subsection*{七层研究链条}
《思想纲领 V1.1》从物理规律出发，依次讨论技术演化、稀缺性变化、经济价值重构、社会制度调整、个人生存战略和精神人生价值\cite{principles}。

从人性需求寻找价值，从供给与稀缺性分析回报，从利益结构研究制度，最后回到个人能做什么。每一层的判断需要相应材料，不能从机器可扩展直接跳到某个人必然失去职业。

\subsection*{七条基本原则}
\begin{enumerate}
\item \textbf{技术演化。}关注机器突破生物智能约束的结构性优势，并检验具体能力与成本。
\item \textbf{需求。}从人需要什么出发，判断成果对谁有作用。
\item \textbf{稀缺性。}持续检查能力和资源的供应、替代与议价条件。
\item \textbf{制度演化。}研究教育、科研、职业、产权与分配怎样回应生产变化。
\item \textbf{双路线生存。}加入AI相关生产与研究，或进入当前较难替代的价值环节。
\item \textbf{精神价值。}重视意义、情绪、信念、真实关系与自主性。
\item \textbf{生存与选择权。}保障基本生活，积累持续调整方向的条件。
\end{enumerate}

\subsection*{生活优先，预案兜底}
AGI预案主张同时认真生活与轻量准备。原稿用“90\%生活成长、10\%风险预备”表达取舍\cite{agi}。这里保留为个人精力安排的提醒，不作为精确最优比例，更不作为投资配置。

承认风险，持续更新认知；同时学习、工作、经营关系和体验生活。光恐惧没有用，也不应让尚未发生的情景吞没当下。

\section{职业选择与科研人的未来}
\subsection*{两条路线可以结合}
纲领提出“加入AI阵营，或进入AI最难替代的阵营”。第一条包括AI研究、应用、智能体、自动化，以及用AI开展科研、工程与产品；第二条从现实需求出发，考察物理资源、实际责任、长期信任、资源控制和共同体\cite{principles}。

这些方向是持续检验的对象，不提供永久安全保证。可以用AI扩大能力，同时参与实际需求与稀缺环节，减少对单一能力的依赖。

\subsection*{数学人的判断框架}
纲领区分数学的内在科学价值与数学劳动的市场稀缺性。常规计算和证明、充分形式化的问题、标准算法、部分建模与重复推理，可能受到影响。这里只保留情景判断，实际案例留在3.5继续收集。

材料提出的转型研究方向包括AI与智能体、工业能源优化、实验与计算、机器人与复杂系统、可靠性与安全关键系统。是否投入，需要结合基础、兴趣、训练环境和真实机会进一步判断。

\subsection*{用成长质量选择导师、项目与环境}
熊辉强调方向感与持续成长。选择导师和项目时，需要关注研究问题、反馈、跨领域合作，以及是否有接触真实问题的机会。短期报酬与平台名称不能独自代替训练质量。

报告建议在实践中形成隐性经验，完成闭环后继续问：“有没有更好的方式？”这可以成为研究习惯：检查能否改进假设、表达、验证或使用方式，逐步培养创新意识。

学习抓主干，并在做事时补细节。让AI辅助查找知识缺口，再通过自己的练习与解释确认掌握。报告关于成长的战略提醒应落实为反馈和修正，不以高薪作为每个人的价值标准。

\section{个人财富与风险管理}
\subsection*{先保障生存，再争取发展}
思想纲领提出保障基本收入与生活条件，建立与自身情况相称的储备，控制不可承受的债务和职业风险，减少对单一技能、行业与收入的依赖\cite{principles}。

储备服务于调整能力，不要求为某个未经验证的未来放弃全部生活。经济安排需要结合实际支出、家庭责任与可承受损失，本稿不设统一金额，也不作具体资产推荐。

\subsection*{两类风险分别准备}
AGI材料区分技术失控与社会分配风险。前者涉及权限、运行与基础设施，后者涉及工作、收入和资源归属。个人可以维护基本生活、联络、资料与转换选择，宏观问题仍需公共制度回应。

不押注单一结局，不购买无法承担的高价末日装备，不因网络预测仓促辞职或脱离现实生活。具体应急行动放在4.6，与可核实情况和当地安排一起判断。

\section{领导力、影响力与社群}
\subsection*{从将兵到将将}
报告把组织多个超级个体作为进一步的培养目标。组织者需要识别能力，提供共同方向，连接伙伴与资源，协调任务和责任，允许成员自主判断\cite{xiong}。

清楚愿景需要持续兑现。成员为什么加入，获得什么支持，怎样表达不同意见，决定如何作出，都应形成明确安排。影响力应与实际帮助和长期信誉联系起来。

\subsection*{连接、吸引、共生}
“创新引力场”的思路是让有能力的人通过平台连接，把孤立资源变成持续合作。平台需要不断提供机会、前沿信息、伙伴与反馈；参与者成长又会吸引新的资源。

“系统成就他人”可以落实到社群：让成员参与提问、资料整理、讨论与实际项目，建立可以重复使用的工作方式。组织成功应包含成员获得能力与机会，不能只看组织者的曝光。

\subsection*{做平台、建系统、定标准、创品牌}
四者可以成为共同体的建设方向：连接人和问题，形成持续协作，说明成果质量与责任，凭长期兑现积累信任。报告强调的“人物”作用，也可以在研究、教学和公共知识中体现。

管理、沟通与同理心需要训练。强个体之间的合作，需要理解他人目标，处理分歧，让不同专长能够配合。平台应扩大成员的能力与联系，而不是让机会全部依附于一个人。

\section{情绪价值、精神力量与人际关系}
思想纲领把精神需求、情绪稳定、真实关系、信念、希望和自主性列为重要主题\cite{principles}。这些需要既存在于个人生活，也存在于社群和共同事业。

用户强调，危机意识应帮助行动，持续恐惧本身不能改善处境。可以承认失落，给担忧安排有限的观察与准备，再把时间放回学习、健康、关系和实际项目。

AGI预案强调不献祭当下。若最坏情景没有发生，生活没有被荒废；若外部环境恶化，已有关系、能力和共同经历也能成为支持。预案只能增加准备，不能保证任何情景下都安全。

真实关系需要长期相处与承担。共同体可以提供归属与合作机会，也应允许不同生活目标。关于具体心理和人际方法，后续有材料再补。

\section{个人行动与长期规划}
\subsection*{把学习、实践与反馈接起来}
沿着思想纲领，持续学习AI、跟踪能力与产业变化，积累可迁移能力、项目经验和社会关系。根据新材料调整方向，同时保留基本生活和选择空间。

报告给出的行动方法是寓教于做、知行合一。走进真实问题，完成一轮交付，接收具体反馈，再问能否做得更好。讨论进一步强调需求访谈、尽早交付原型和持续迭代\cite{xiong,discussion}。

学习可以从知识框架开始，诊断缺口，补充必要内容，再检验是否能够独立解释和使用。选择环境时，检查是否存在专业反馈、合作伙伴与实际机会。没有这些条件，忙碌未必形成持续成长。

\subsection*{极端风险的分级预案}
原AGI材料提出黄、橙、红、黑四级，用来把担忧转成分阶段行动\cite{agi}。这是个人情景框架，尚未验证预警精度，也不是官方等级。AI限流、监管变化或一次系统故障，不能单独确认灾难；应核实事实，结合实际服务状况与当地应急信息。

\begin{center}\small
\begin{tabularx}{\textwidth}{P{0.13\textwidth}YY}
\toprule 等级 & 原稿讨论的状态 & 个人行动重心\\\midrule
黄 & 风险关注增加，基本生活正常 & 维持日常，核实消息，检查基础准备与联络\\
橙 & 多项可信异常，服务出现不稳定 & 更新计划，补齐合理缺口，避免抢购与盲动\\
红 & 本地供水供电、交通或秩序受到严重影响 & 依据当地指引决定避险，优先生活安全与联络\\
黑 & 极端长期中断情景 & 依靠可信合作与可用公共支持，重新评估长期安排\\
\bottomrule
\end{tabularx}
\end{center}

原稿包含不同周数和月数的储备，以及离线、撤离等设想。本稿不把它们改写成适合所有人的固定指令。储备要结合资源、轮换与实际需要；撤离取决于具体威胁、路线和指引，断网也可能使人失去求助信息。

\subsection*{行动边界}
保留原材料的几条底线：不把个人对抗AGI作为方案，不自制EMP、爆炸物等危险装置，不盲目购买高价末日产品，不传播未经核实的灾难消息，不把全部人生押在一个判断上。

预案的目标是减少临时慌乱，保持可执行的安排。能力跃升、隐蔽失效与社会动荡都属于需要研究的情景，不能仅凭想象确定发生时间。日常仍以认真生活、成长与实际合作为主。

\backmatter
\renewcommand{\bibname}{本稿采用的素材}
\begin{thebibliography}{9}
\small\raggedright
\bibitem{project}《AI时代生存指南知识库·项目介绍(2)》PPT。采用第2、3页目录，第4至6页前言与行动原则。
\bibitem{principles}《〈AI时代生存指南〉思想纲领 V1.1》PPT。采用七层链条、需求与稀缺性、数学人专题、双路线、生存原则与精神价值等内容。假说按假说整理。
\bibitem{agi}《AGI极端风险·个人终极认知与生存预案体系》PPT。采用风险与意识分离、极端情景、双重危机、生活与准备取舍、分级预案和行动边界。对其中必然性、概率和具体避险指令不作科学或通用安全背书。
\bibitem{xiong}熊辉，《熊辉副校长报告原文》整理稿，用户于2026年10月27日提供。采用人才培养、成长方式、人机团队、隐性经验、知识转移、创新生态与制度思考。未将材料中无原始出处的薪资、产能与估值数字用作一般结论。
\bibitem{discussion}《AI时代高校教育与创业相关讨论》带时间戳转写，标注2026年10月27日10:23至11:25。采用教学尝试、产业合作、需求验证和早期创业反馈。转写可能存在识别错误。
\bibitem{summary}同场讨论的智能纪要摘要版。与转写对照使用，不视为独立证据。摘要中的确定表述按原发言中的保留意见处理。
\end{thebibliography}
\end{document}
